\documentclass[10pt]{mypackage}

\usepackage[uprightgreeks,light,noDcommand]{kpfonts}
\renewcommand{\coloneq}{\coloneqq}
\renewcommand{\eqcolon}{\eqqcolon}
%\renewcommand{\mathbb}{\mathds}
%\usepackage{homework}
%\usepackage{notes}

%\usepackage[ backend=bibtex, style = alphabetic, sorting=ynt ]{biblatex}
%\addbibresource{exam_prep_syllabus.bib}

\usepackage{parskip}

%\fancyhf{}
%\fancyhead[R]{Avinash Iyer}
%\fancyhead[L]{}
%\fancyfoot[C]{\thepage}
\title{Oral Exam Syllabus}
\date{October 2026}
\author{Avinash Iyer}

\setcounter{secnumdepth}{0}

\begin{document}
\RaggedRight
\maketitle
The oral exam for Avinash Iyer will be administered by Professors Ben Hayes and [either Mikhail Ershov or Thomas Koberda] on December 4, 2026 focused on Functional Analysis and Geometric Group Theory. The exam will consist of questions on both topics.

\textbf{Functional Analysis Topics:}
\begin{itemize}
  \item Locally Convex Spaces, topics covering Chapter IV \S 1--\S 3 of Conway's \textit{A Course in Functional Analysis}
    \begin{itemize}
      \item Definitions and examples
      \item Metrizability and Normability
      \item Geometric Hahn--Banach
    \end{itemize}
  \item Weak topologies and extremal structure, topics covering Chapter V \S 1--\S 10 \S 12--\S 13 of Conway's \textit{A Course in Functional Analysis}
    \begin{itemize}
      \item Dual spaces of locally convex topologies
      \item Banach--Alaoglu Theorem
      \item Reflexivity, Separability, and Metrizability
      \item Stone--\v{C}ech compactification
      \item Krein--Milman theorem
      \item Schauder Fixed-Point Theorem, Ryll--Nardzewski Fixed-Point Theorem
      \item Krein--Smulian Theorem
    \end{itemize}
  \item $C^{\ast}$-algebras and von Neumann algebras, topics covering Chapter VIII \S 1--\S 3, \S 5 of Conway's \textit{A Course in Functional Analysis} and Sections 4.6 and 4.7 of Pedersen's \textit{Analysis Now}
    \begin{itemize}
      \item Properties and examples
      \item Gelfand Transform, Gelfand--Naimark Theorem, Continuous Functional Calculus
      \item Ordering properties
      \item GNS Construction
      \item von Neumann double commutant theorem
      \item $\sigma$-topologies and the predual of a von Neumann algebra
    \end{itemize}
  \item Spectral theory for normal operators, topics covering Chapter IX \S 1--\S 8 of Conway's \textit{A Course in Functional Analysis} and Sections 4.4--4.6 of Pedersen's \textit{Analysis Now}
    \begin{itemize}
      \item spectral measures and projections
      \item Borel functional calculus
      \item the characterization of abelian von Neumann algebras
    \end{itemize}
\end{itemize}
\textbf{Geometric Group Theory:}
\begin{itemize}
  \item Relations between groups and their action on graphs, topics covering Chapters 3 and 4 of Löh's \textit{Geometric Group Theory}
    \begin{itemize}
      \item Cayley graphs
      \item actions of groups on trees; proof of the Nielsen--Schreier Theorem
    \end{itemize}
  \item Quasi-isometry, topics covering Chapter 5 and 6 of Löh's \textit{Geometric Group Theory} and sections 11.12--11.13 of Ceccherini-Silberstein and d'Adderio's \textit{Topics in Groups and Geometry}
    \begin{itemize}
      \item definition of quasi-isometry of metric spaces and groups, quasi-geodesics;
      \item \v{S}varc--Milnor lemma and its topological formulation (statement and applications)
      \item dynamic criterion for quasi-isometry (set-theoretic coupling and topological coupling)
      \item asymptotic cones
      \item notion of growth rate
      \item quasi-isometry invariance of growth rate
    \end{itemize}
  \item Amenability, topics covering Chapter 9 of Löh's \textit{Geometric Group Theory} and other sources
    \begin{itemize}
      \item definition via means and Følner Sequences, elementary inheritance
      \item proofs of sufficient criteria (polynomial growth, solvability)
      \item equivalent criteria (paradoxicality, almost-invariant vectors, weak containment)
      \item quasi-isometry invariance of amenability
    \end{itemize}
  \item Soficity, topics covering content from Pestov's \textit{Hyperlinear and Sofic Groups}, Pestov and Kwiatowska's \textit{An Introduction to Hyperlinear and Sofic Groups}, and Ceccherini-Silberstein and Coornaert's \textit{Cellular Automata and Groups}
    \begin{itemize}
      \item almost-homomorphism into symmetric groups;
      \item embedding into metric ultraproduct (and sofic approximations);
      \item initial subamenability of Cayley graph;
      \item inheritance properties and prominent examples (amenable groups, residually finite groups)
    \end{itemize}
\end{itemize}
\end{document}
